\chapter{Chikungunya}

\section*{Definition and Overview}
Chikungunya is a viral illness spread by the bite of infected mosquitoes, mainly the daytime-biting \textit{Aedes} species. The name comes from a word in the Kimakonde language meaning \textit{to become contorted}, reflecting the severe joint pain that is the hallmark of the disease. Chikungunya typically causes the sudden onset of high fever and painful joints, often accompanied by a rash. It is rarely fatal, but the joint pain can persist for months and, in some people, for years, causing substantial disability. There is no specific antiviral treatment, so management focuses on symptom relief and rehabilitation.

\section*{Epidemiology}
Chikungunya occurs in tropical and subtropical regions of Africa, Asia, the Indian Ocean islands, the Pacific, and the Americas. Outbreaks are often explosive, and millions of cases have occurred during large epidemics in recent decades. The disease is also seen in travellers returning from affected areas, occasionally causing locally acquired clusters in regions where the \textit{Aedes} mosquito is present. Transmission follows the distribution of the mosquito vectors and tends to be seasonal, with peaks after rains. Anyone can be infected, but severe disease is more likely in newborns, older adults, and people with underlying medical conditions.

\section*{Aetiology and Pathophysiology}
\medskip
\textbf{Mechanism of disease.} Chikungunya is caused by an alphavirus transmitted when a mosquito takes a blood meal from an infected person and later bites another person:
\begin{itemize}
    \item \textbf{Entry and spread:} after the virus enters the skin through the mosquito bite, it spreads through the bloodstream and infects cells in the joints, muscles, and skin
    \item \textbf{Inflammation:} infection provokes a strong inflammatory response that is responsible for the fever, joint swelling, and pain
    \item \textbf{Viral clearance:} the immune system normally clears the virus within about one to two weeks
    \item \textbf{Persistent joint symptoms:} in some people, joint inflammation continues after the virus has gone, a process thought to involve ongoing immune and inflammatory responses
\end{itemize}

\section*{Clinical Features}
\begin{itemize}
    \item \textbf{Fever:} sudden onset of high fever, often with chills
    \item \textbf{Joint pain:} severe, typically affecting the hands, wrists, ankles, and knees, often on both sides
    \item \textbf{Joint swelling:} redness, swelling, and stiffness of the affected joints
    \item Muscle aches and headache
    \item A red, flat rash over the trunk and limbs in many patients
    \item Fatigue and weakness that may last for several weeks
    \item Less commonly, nausea, vomiting, and pain behind the eyes
\end{itemize}

\section*{Differential Diagnosis}
\begin{itemize}
    \item Dengue fever -- a similar mosquito-borne illness with fever, aches, and sometimes bleeding
    \item Zika virus infection -- a milder illness with fever, rash, and joint pain
    \item Other arboviruses, such as Ross River virus and Barmah Forest virus
    \item Malaria, particularly in a returning traveller
    \item Measles or rubella presenting with a febrile rash
    \item Early rheumatoid arthritis when joint symptoms persist
\end{itemize}

\section*{When to Seek Medical Care}
Seek medical assessment if you develop a fever and joint pain after travelling to a region where chikungunya occurs, or if symptoms are severe, unusual, or prolonged.

\medskip
\textbf{Call 000 (or local emergency services) for:}
\begin{itemize}
    \item Persistent high fever with confusion or drowsiness
    \item Difficulty breathing or chest pain
    \item Inability to drink fluids or signs of dehydration
    \item Symptoms in a newborn, an older person, or someone with a chronic illness
    \item Any unusual bleeding or bruising
\end{itemize}

\section*{Diagnosis and Investigations}
\begin{itemize}
    \item \textbf{History:} recent travel to an affected area and a history of mosquito exposure
    \item \textbf{PCR testing:} within the first week of illness, blood tests can detect the viral genetic material directly
    \item \textbf{Serology:} antibody blood tests (IgM and IgG) confirm the diagnosis from around the end of the first week
    \item \textbf{Full blood count:} a low white cell and platelet count is common but not specific
    \item \textbf{Exclusion of other causes:} tests for dengue, malaria, and other viruses are often needed, as the illnesses overlap
    \item \textbf{Assessment of complications:} further investigations if neurological, cardiac, or other complications are suspected
\end{itemize}

\section*{Management}
\begin{itemize}
    \item \textbf{Rest and fluids:} important for recovery and to prevent dehydration from fever
    \item \textbf{Pain relief:} paracetamol for fever and pain, taken in the recommended dose
    \item \textbf{Anti-inflammatory medicines:} such as ibuprofen to relieve joint swelling and pain once dengue has been excluded
    \item \textbf{Physiotherapy:} gentle exercise and stretching once the acute fever settles, to maintain joint mobility
    \item \textbf{Monitoring:} review by a doctor if joint pain or fatigue persists beyond a few weeks
    \item \textbf{No specific antiviral:} treatment is supportive, as the immune system clears the virus
\end{itemize}

\section*{Complications}
\begin{itemize}
    \item Persistent or chronic joint pain lasting months or years, especially in older adults
    \item Chronic fatigue and reduced quality of life
    \item Rare neurological problems, including inflammation of the brain, nerves, or spinal cord
    \item Severe illness in newborns infected around the time of birth
    \item Deterioration of underlying conditions such as heart disease or diabetes
\end{itemize}

\section*{Prognosis and Outcomes}
Most people recover fully within one to two weeks. However, joint pain persists for weeks to months in a substantial number of patients, and a minority continue to experience pain and stiffness for a year or more, which can affect work and daily life. Death is uncommon and occurs mainly in newborns and frail older adults. Because there is no specific treatment, recovery depends on the immune system, effective symptom relief, and gradual rehabilitation. People with persistent joint symptoms should be followed up, as they may benefit from rheumatology review.

\section*{Prevention and Self-Care}
\begin{itemize}
    \item Prevent mosquito bites with insect repellent containing DEET applied to exposed skin
    \item Wear long sleeves and long trousers, especially during the day when \textit{Aedes} mosquitoes are most active
    \item Use insecticide-treated bed nets and window screens in affected areas
    \item Remove standing water around the home where mosquitoes breed
    \item Travellers to outbreak areas should take these precautions seriously
    \item A vaccine is available in some countries for people at risk; discuss this with a doctor before travel
\end{itemize}

\section*{Sources and Further Reading}
The following reputable sources provide further information on chikungunya:
\begin{itemize}
    \item World Health Organization. \textit{Chikungunya fact sheets and outbreak information.} https://www.who.int/
    \item Centers for Disease Control and Prevention. \textit{Chikungunya resources for health professionals.} https://www.cdc.gov/
    \item NHS. \textit{Health A--Z: information on travel-related illnesses.} https://www.nhs.uk/conditions/
    \item Mayo Clinic. \textit{Diseases and conditions library.} https://www.mayoclinic.org/
    \item MSD Manual. \textit{Professional edition: viral infections.} https://www.msdmanuals.com/
    \item MedlinePlus. \textit{Health topics on infectious diseases.} https://medlineplus.gov/
\end{itemize}
